\chapter{Lập lịch và cân bằng dây chuyền}
\chapterlead{Một phép biểu diễn cho bài toán thời gian chỉ thực sự rõ khi trả lời trực tiếp
ba câu hỏi: quyết định nằm ở mốc nào, điều gì phải liên tục giữa các mốc, và
cận mục tiêu được thay đổi ra sao trong quá trình giải.}

\index{scheduling}\index{non-preemption}\index{Nurse Rostering}
\index{lazy clause generation}\index{Dynamic Discretization Discovery}

\flowdiagram{Miền thời gian}{Biến quyết định}{Liên tục và tài nguyên}
  {Hàm mục tiêu}{Chứng nhận}

\section{Từ SAT theo mốc thời gian đến lazy clause generation}

Lập lịch là một trong những miền sớm cho thấy SAT có thể giải các bài toán
ngoài lôgic thuần túy. Các thí nghiệm SAT về lập lịch từ thập niên 1990 đã
dùng quyết định theo mốc thời gian và mệnh đề xung đột để so sánh với phương
pháp tìm kiếm chuyên dụng \parencite{crawford1994scheduling}. Từ đó, văn liệu
phát triển theo ba hướng:
\begin{enumerate}
  \item biên dịch toàn bộ ràng buộc gán, thời gian và tài nguyên sang SAT;
  \item dùng MaxSAT/PB để biểu diễn mức phạt, ưu tiên hoặc số ca vi phạm;
  \item tích hợp lan truyền miền với học mệnh đề, điển hình là
        \emph{lazy clause generation} (sinh mệnh đề lười, LCG).
\end{enumerate}

LCG giữ biến miền và bộ lan truyền của CP, đồng thời gắn mỗi phép loại giá trị
với một giải thích dạng mệnh đề; phân tích xung đột sau đó học trên các giải
thích này. Trên RCPSP/max, kiến trúc đó kết hợp lan truyền lập lịch toàn cục
với học mệnh đề \parencite{schutt2013lcg}. Ở một hướng khác, MaxSAT bộ phận có
trọng số cho phép tách độ phủ/tính hợp lệ thành các mệnh đề cứng và ưu tiên
trong lập lịch nhân viên thành các mệnh đề mềm
\parencite{demirovic2019staff}.

\begin{center}
\captionof{table}{Các mô hình giải chính xác cho bài toán thời gian.}
\label{tab:scheduling-paradigms}
\begin{tabularx}{\textwidth}{@{}>{\bfseries\sffamily}p{.18\textwidth}YYY@{}}
\toprule
Mô hình & Trạng thái thời gian & Cơ chế tài nguyên & Cận chủ đạo\\
\midrule
SAT theo mốc thời gian & Mốc rời rạc & Xung đột/ràng buộc đếm & UNSAT theo khung thời gian\\
SAT thứ tự & Ngưỡng thời gian & Quan hệ + cận dịch chuyển & Học mệnh đề\\
MaxSAT/PB & Như SAT + literal chi phí & Cứng/mềm hoặc PB & Lõi/cận chi phí\\
SMT & Biến nguyên/thực & Phép tuyển + lý thuyết & Lan truyền lý thuyết\\
CP/LCG & Miền và khoảng & Bộ lan truyền toàn cục & Lọc miền + mệnh đề\\
MIP & Biến nhị phân/liên tục & Bất đẳng thức tuyến tính & Nới lỏng quy hoạch tuyến tính\\
\bottomrule
\end{tabularx}
\end{center}

Vì vậy, một phép biểu diễn lập lịch được đánh giá trên cả kích thước, sức lan
truyền và biểu diễn đối chứng tương ứng: SAT theo mốc thời gian so với MIP mở
rộng theo thời gian; SAT thứ tự so với lôgic sai phân/SMT; bộ đếm chuỗi dùng
chung so với bộ lan truyền \textsc{Sequence}; và hàm mục tiêu MaxSAT so với PB
hay hàm mục tiêu tuyến tính.

\section{Kiến trúc chung của mô hình thời gian}

Các bài toán trong chương có bề mặt khác nhau: xếp tác vụ không ngắt
trên các máy đồng nhất, phân ca điều dưỡng, lập lịch cuộc họp B2B với cân bằng thời
gian chờ, cân bằng dây chuyền theo đỉnh công suất và điều chỉnh lịch tàu. Tuy
nhiên, phép biểu diễn của chúng đều có thể tách thành bốn lớp:
\begin{enumerate}
  \item \textbf{gán}: tác vụ, người hoặc chuyến tàu sử dụng tài nguyên nào;
  \item \textbf{thời gian/thứ tự}: thời điểm bắt đầu, kết thúc hoặc thứ tự trước--sau;
  \item \textbf{liên tục/chuỗi}: các quyết định cục bộ ghép thành một
        khoảng hoặc một chuỗi hợp lệ;
  \item \textbf{hàm mục tiêu}: thời điểm hoàn tất toàn lịch, đỉnh công suất,
        độ trễ hay số vi phạm.
\end{enumerate}

Phân tách kiến trúc theo lớp không chỉ giúp trình bày mạch lạc, mà còn cho phép thay thế một mô-đun mà không
đổi ngữ nghĩa của các mô-đun còn lại. Chẳng hạn, có thể giữ nguyên biến gán
máy và thay phép biểu diễn cho tính không ngắt; hoặc giữ nguyên mô hình khả thi
và thay vòng lặp tối ưu bằng \SAT tăng dần hay \MaxSAT.

\begin{designrule}{Nguyên tắc thiết kế}
Mỗi biến phụ phải thuộc về một lớp rõ ràng. Nếu một biến vừa được dùng để suy
thời điểm, vừa được diễn giải như trạng thái hoạt động, cần có liên kết hai
chiều hoặc một phép giải mã chứng minh được.
\end{designrule}

\section{Lập lịch không ngắt trên các tài nguyên đồng nhất}

\subsection{Mô hình}

Cho tập tác vụ \(J=\set{1,\ldots,n}\) và \(m\) tài nguyên đồng nhất. Tác vụ
\(i\) có thời điểm sẵn sàng \(r_i\), thời lượng \(e_i\) và hạn hoàn tất
\(d_i\). Thời điểm bắt đầu thuộc miền
\[
  T_i=\set{r_i,r_i+1,\ldots,d_i-e_i}.
\]
Nếu \(s_i\in T_i\), tác vụ chiếm tài nguyên trong nửa khoảng
\([s_i,s_i+e_i)\). Hai tác vụ được gán cùng tài nguyên không được có hai khoảng
giao nhau.

Một mô hình gọn dùng hai họ biến:
\[
  y_{ij}=1 \iff i\text{ được gán cho tài nguyên }j,
  \qquad
  z_{it}=1 \iff i\text{ hoạt động tại }t .
\]
Với mỗi tác vụ, các biến \(y_{ij}\) thỏa
\[
  \ExactlyOne(y_{i1},\ldots,y_{im}).
\]
Xung đột tài nguyên tại mốc \(t\) được loại bởi
\[
  \neg z_{it}\lor \neg y_{ij}\lor
  \neg z_{i't}\lor \neg y_{i'j},
  \qquad i<i'.
\]
Mệnh đề này thể hiện xung đột có điều kiện: mệnh đề chỉ kích hoạt khi hai tác vụ cùng
chọn tài nguyên \(j\).

\subsection{Từ ràng buộc đếm đến một khối liên tục}

Điều kiện
\[
  \sum_t z_{it}=e_i
\]
chỉ bảo đảm đủ số mốc hoạt động; cơ chế này không ngăn các mốc bị phân tách thành nhiều
đoạn. \emph{Non-preemption} (tính không ngắt: tác vụ phải chạy liên tục từ
lúc bắt đầu đến lúc kết thúc) đòi hỏi mạnh hơn:
\[
  z_{it}=1 \iff s_i\le t<s_i+e_i .
\]

Direct Encoding có thể tạo một biến \(b_{is}\) cho từng thời điểm bắt đầu
\(s\in T_i\). Trước hết phải đóng lựa chọn thời điểm bắt đầu bằng
\[
  \ExactlyOne\bigl(b_{is}:s\in T_i\bigr).
\]
Trong \CNF, phần ít nhất một là mệnh đề
\(\bigvee_{s\in T_i}b_{is}\); phần nhiều nhất một được sinh bằng biểu diễn từng
cặp, Sequential Counter hoặc một mô-đun AMO khác tùy kích thước \(T_i\). Sau đó
dùng
\[
  b_{is}\rightarrow z_{it}
  \quad\text{với }s\le t<s+e_i,
  \qquad
  b_{is}\rightarrow \neg z_{it}
  \quad\text{ở ngoài khoảng.}
\]
Cách này dễ hiểu nhưng lặp lại nhiều literal cho các thời điểm bắt đầu lân cận.
Phép biểu diễn gọn thay phần lặp bằng các biến tiền tố/hậu tố có ngữ nghĩa ``tất
cả bằng 0'' hoặc ``tất cả bằng 1''. Hai thời điểm bắt đầu kế tiếp chia sẻ gần
như toàn bộ cấu trúc; khác biệt của chúng chỉ nằm ở hai biên của khối.

\begin{proposition}[Đúng đắn của phép biểu diễn theo khối]
\label{prop:nonpreemptive-block}
Giả sử mô-đun khối bảo đảm: (i) tồn tại đúng một vị trí bắt đầu; (ii) từ vị
trí đó có đúng \(e_i\) giá trị liên tiếp bằng \(1\); và (iii) mọi vị trí trước
và sau khối đều bằng \(0\). Khi đó phép gán cho các biến \(z_{it}\) tương ứng
duy nhất với một khoảng xử lý không ngắt hợp lệ của tác vụ \(i\).
\end{proposition}

\begin{proof}
Từ (i), chọn vị trí bắt đầu \(s_i\). Điều kiện (ii) cho
\(z_{it}=1\) trên \([s_i,s_i+e_i)\); điều kiện (iii) loại mọi giá trị \(1\)
khác. Chiều ngược lại, từ một khoảng hợp lệ, đặt biến bắt đầu ở đầu khoảng và
gán các biến tiền tố/hậu tố theo định nghĩa của chúng. Mọi mệnh đề của mô-đun
đều được thỏa.
\end{proof}

\begin{figure}[tb]
\centering
\begin{tikzpicture}[satfig,x=9mm,y=8mm]
  \draw[->,Rule] (0,0)--(9,0) node[right,satlabel]{thời gian};
  \foreach \t in {0,...,8}{
    \draw[Rule] (\t,.12)--(\t,-.12);
    \node[below,satlabel] at (\t,-.15) {\(\t\)};
  }
  \node[left,satlabel] at (0,1.05) {hợp lệ};
  \draw[Blue,fill=PaleBlue,thick] (2,.55) rectangle (5,1.45);
  \node at (3.5,1) {tác vụ \(i\), \(e_i=3\)};
  \node[left,satlabel] at (0,2.55) {không hợp lệ};
  \draw[Gold,fill=PaleGold,thick] (1,2.1) rectangle (2,3);
  \draw[Gold,fill=PaleGold,thick] (4,2.1) rectangle (6,3);
  \node[satlabel,text=Gold] at (3.5,3.35)
    {đủ \(3\) mốc nhưng bị ngắt};
\end{tikzpicture}
\caption{Ràng buộc đếm \(\sum_tz_{it}=3\) không phân biệt một khối liên tục với
ba mốc rời; tính không ngắt phải biểu diễn vị trí bắt đầu và hai biên khối.}
\label{fig:nonpreemptive-block}
\end{figure}

\begin{workedexample}{Ba tác vụ trên hai máy}
Cho \(J_1=(r=0,e=3,d=6)\), \(J_2=(0,2,5)\), \(J_3=(1,2,6)\).
Một lịch hợp lệ là
\[
  M_1:\ J_1[0,3)\ \ J_3[3,5),
  \qquad
  M_2:\ J_2[0,2).
\]
Với \(t=1\), \(z_{1,1}=z_{2,1}=1\) nhưng hai tác vụ dùng máy khác nhau nên
mệnh đề xung đột có điều kiện vẫn thỏa. Nếu đổi \(y_{2,2}\) thành \(y_{2,1}\),
mệnh đề tại \(t=1,j=1\) trở thành
\(
\neg z_{1,1}\lor\neg y_{1,1}\lor
\neg z_{2,1}\lor\neg y_{2,1}
\)
và bị vi phạm. Ví dụ tách rõ hai lớp: \(z\) nói khi nào tác vụ chạy, \(y\)
nói tác vụ chạy ở vị trí nào.
\end{workedexample}

Kết quả của \textcite{kieu2025scheduling} cho thấy lợi ích chính của phép mã
hóa gọn nằm ở việc tránh lặp cấu trúc liên tục. Cấu hình kết hợp biểu diễn theo
khối với phá vỡ đối xứng và một phép biểu diễn giả Boolean thích hợp đạt tổng thời
gian thấp nhất trong nhiều nhóm bài toán của nghiên cứu đó. Điều cần rút ra cho
thiết kế là: \emph{tính không ngắt nên là một mô-đun độc lập}, thay vì được
ngầm suy từ một ràng buộc đếm.

\section{Ràng buộc chuỗi trong lập lịch điều dưỡng}

\subsection{Cửa sổ trượt}

Với dãy Boolean \(\Omega=(x_1,\ldots,x_n)\), một ràng buộc chuỗi cận dưới yêu
cầu mọi cửa sổ độ rộng \(w\) chứa ít nhất \(k\) giá trị đúng:
\[
  \operatorname{ALSC}(\Omega,w,k)=
  \bigwedge_{i=0}^{n-w}
  \left(\sum_{j=i+1}^{i+w}x_j\ge k\right).
\]
Nếu biểu diễn từng cửa sổ bằng một phép biểu diễn ràng buộc đếm độc lập, hai cửa sổ
kề nhau chia sẻ \(w-1\) biến gốc nhưng không chia sẻ biến phụ. Cấu trúc chuỗi
vì thế đặt ra bài toán \emph{tái sử dụng biểu thức con}.

Chia một cửa sổ thành hai phần \(A\) và \(B\). Ta có
\[
  \sum_{x\in A\cup B}x\ge k
  \iff
  \bigwedge_{q=1}^{k}
  \left[
    \left(\sum_{x\in A}x\ge q\right)
    \lor
    \left(\sum_{x\in B}x\ge k-q+1\right)
  \right].
\]
Mỗi mệnh đề ngưỡng trong ngoặc được biểu diễn bởi một thanh ghi chia sẻ.
Khi cửa sổ trượt sang phải, phần lớn thanh ghi vẫn còn nguyên.

\begin{example}
Xét \(x_3+x_4+x_5+x_6\ge2\), chia thành hai khối
\(A=(x_3,x_4)\) và \(B=(x_5,x_6)\). Nếu
\(R_{A,q}\) mang nghĩa ``\(A\) có ít nhất \(q\) giá trị đúng'', bộ nối chỉ
cần
\[
  (R_{A,1}\lor R_{B,2})
  \land
  (R_{A,2}\lor R_{B,1}).
\]
Các thanh ghi của \(B\) tiếp tục được dùng trong cửa sổ kế tiếp.
\end{example}

\begin{figure}[tb]
\centering
\begin{tikzpicture}[satfig,x=12mm,y=7mm]
  \foreach \d/\v in {1/1,2/0,3/1,4/0,5/1,6/0,7/1}{
    \ifnum\v=1
      \node[satbox,minimum width=10mm,fill=PaleTeal] (d\d) at (\d,0) {\(\v\)};
    \else
      \node[satbox,minimum width=10mm,fill=SoftGray] (d\d) at (\d,0) {\(\v\)};
    \fi
    \node[above,satlabel] at (d\d.north) {ngày \(\d\)};
  }
  \draw[Blue,very thick] (.55,-.75)--(4.45,-.75)
    node[midway,below,satlabel,text=Blue]{\(W_1\): tổng \(2\)};
  \draw[Teal,very thick] (1.55,-1.55)--(5.45,-1.55)
    node[midway,below,satlabel,text=Teal]{\(W_2\): tổng \(2\)};
  \draw[Gold,very thick] (2.55,-2.35)--(6.45,-2.35)
    node[midway,below,satlabel,text=Gold]{\(W_3\): tổng \(2\)};
\end{tikzpicture}
\caption{Ba cửa sổ ALSC \((w=4,k=2)\) trên bảy ngày. Mỗi cửa sổ kề nhau chỉ
thay một ngày, nên các ngưỡng ở phần giao có thể dùng lại.}
\label{fig:nurse-alsc}
\end{figure}

\subsection{Đưa vào mô hình phân ca}

Đặt \(x_{ids}=1\) khi điều dưỡng \(i\) làm loại ca \(s\) ở ngày \(d\). Mỗi cặp
\((i,d)\) nhận đúng một trạng thái. Những quy tắc như ``trong mọi \(w\) ngày có
ít nhất \(k\) ngày nghỉ'' trở thành ALSC; quy tắc at-most được đối ngẫu:
\[
  \sum_{j=1}^{w}x_j\le q
  \iff
  \sum_{j=1}^{w}(1-x_j)\ge w-q.
\]
Nhờ đó một bản cài đặt có thể dùng cùng một mô-đun cận dưới cho cả hai hướng.

Trong một mô hình phân ca đầy đủ, ràng buộc chuỗi chỉ là một lớp. Lớp độ phủ
bảo đảm số nhân viên mỗi ca; ràng buộc hợp đồng giới hạn tổng giờ và số cuối
tuần; ưu tiên tạo các chi phí mềm. Phép biểu diễn cửa sổ chia sẻ tác động trực
tiếp lên lớp chuỗi, còn MaxSAT có trọng số là một lựa chọn cho hàm mục tiêu của
toàn mô hình \parencite{demirovic2019staff}. Ma trận thiết kế tách các lớp để
đo riêng ảnh hưởng của ALSC, độ phủ và ưu tiên.

\begin{resultbox}{Kết quả tiêu biểu}
Trên 28 bài toán \emph{Nurse Rostering} (lập lịch điều dưỡng), với quy mô đến
120 điều dưỡng và 24 tuần,
ALSCE giải toàn bộ tập thử trong nghiên cứu của
\textcite{truong2025nurse}; bài toán lớn nhất được báo cáo ở mức
\(26{,}01\) giây. Kết quả này minh họa hiệu quả của kỹ thuật chia sẻ thanh ghi giữa
các cửa sổ chồng lấn.
\end{resultbox}

\section{Lập lịch cuộc họp B2B với cân bằng thời gian chờ}

\index{B2B meeting scheduling}\index{idle-time balancing}
\index{SparseSuffix}\index{domain filtering}

\subsection{Bài toán B2B và mục tiêu khoảng thời gian chờ}

Bài toán lập lịch cuộc họp doanh nghiệp (Business-to-Business Meeting Scheduling Problem, B2B) xếp lịch cho tập các cuộc họp đối tác \(M\) vào tập mốc thời gian rời rạc \(T=\set{1,\ldots,h}\) và \(K\) phòng họp/địa điểm dùng chung \parencite{pesant2015b2b,bofill2022b2b}. Mỗi cuộc họp \(m\in M\) được thực hiện bởi đúng hai đại biểu, ký hiệu \(\pi(m)\subseteq P\) với \(\card{\pi(m)}=2\), và có miền mốc khởi tạo \(D_m^0\subseteq T\) tương ứng với phiên làm việc, mốc cố định và mốc cấm của hai đại biểu. Đặt \(E\subseteq M\times M\) là tập các cung thứ tự tiên quyết giữa các cuộc họp. Một lịch xếp \(S: M\to T\) là hợp lệ nếu \(S(m)\in D_m^0\) và thỏa mãn các điều kiện:
\begin{align}
  S(m) &\ne S(m') && \forall m,m'\in M, m\ne m': \pi(m)\cap\pi(m')\ne\emptyset, \label{eq:b2b-participant}\\
  \card{\set{m\in M \suchthat S(m)=t}} &\le K && \forall t\in T, \label{eq:b2b-capacity}\\
  S(i) &< S(j) && \forall (i,j)\in E. \label{eq:b2b-precedence}
\end{align}

Với mỗi đại biểu \(p\in P\), đặt \(M_p = \set{m\in M \suchthat p\in\pi(m)}\) là tập cuộc họp của đại biểu đó. Đặt \(P^\star = \set{p\in P \suchthat \card{M_p}\ge 2}\) là tập các đại biểu có ít nhất hai cuộc họp. Đối với \(p\in P^\star\), mốc bận đầu tiên \(F_p(S)\) và mốc bận cuối cùng \(L_p(S)\) trong lịch \(S\) được xác định bởi:
\[
  F_p(S) = \min_{m\in M_p} S(m), \qquad L_p(S) = \max_{m\in M_p} S(m).
\]
Do các cuộc họp của đại biểu \(p\) không chồng lấn, số mốc thời gian trống nằm hoàn toàn giữa mốc đầu và mốc cuối---gọi là \emph{thời gian chờ nội bộ} (internal idle slots)---được tính bằng:
\begin{equation}
  I_p(S) =
  \begin{cases}
    L_p(S) - F_p(S) + 1 - \card{M_p}, & p\in P^\star,\\
    0, & \card{M_p}\le 1.
  \end{cases}
  \label{eq:b2b-internal-idle}
\end{equation}
Khác với mô hình truyền thống đếm tổng số nhóm nghỉ (\emph{break groups}) \parencite{bofill2022b2b}, chỉ số \(I_p(S)\) phân biệt rõ khoảng chờ dài (như chờ 5 mốc) với khoảng chờ ngắn (chờ 1 mốc). Hàm mục tiêu tối ưu hóa khoảng thời gian chờ nội bộ giữa các đại biểu:
\begin{equation}
  \Delta_I(S) = \max_{p\in P^\star} I_p(S) - \min_{p\in P^\star} I_p(S).
  \label{eq:b2b-idle-range}
\end{equation}
Mục tiêu này hướng tới sự cân bằng công bằng về thời gian chờ giữa các đại biểu, lấy cảm hứng từ các chuẩn cân bằng tải trong bài toán BACP \parencite{chiarandini2012}. Tổng thời gian chờ \(I_\Sigma(S)=\sum_{p\in P}I_p(S)\) được báo cáo như một chỉ số bổ trợ.

\subsection{Tiền xử lý bảo toàn nghiệm và lọc miền}

Trước khi sinh công thức \CNF, thuật toán lọc miền \parencite{nguyen2026compactsatmaxsatencodings} thu hẹp miền mốc của từng cuộc họp bằng cách áp dụng lặp cho đến điểm bất động ba quy tắc lọc:
\begin{enumerate}
  \item \textbf{Cận thứ tự tiên quyết}: Với mỗi quan hệ \((i,j,d)\) trong tập quan hệ \(\mathcal R_f\in\set{\mathcal R_E, \mathcal R_{E^\star}}\) (trong đó \(\mathcal R_E = \set{(i,j,1) \suchthat (i,j)\in E}\) và \(\mathcal R_{E^\star}\) là bao đóng bắc cầu gán khoảng cách đường đi dài nhất \(\ell_{ij}\)), loại \(t\in D_i\) nếu \(t+d > \max D_j\), và loại \(u\in D_j\) nếu \(\min D_i+d > u\).
  \item \textbf{Nhất quán ghép cặp đại biểu}: Với mỗi đại biểu \(p\), dựng đồ thị hai phía giữa \(M_p\) và miền mốc khả thi. Sử dụng thuật toán lọc \textsc{AllDifferent} dựa trên ghép cặp cực đại \parencite{regin1994alldiff} để xóa phân gán cuộc họp--mốc nào không thuộc bất kỳ ghép cặp hoàn hảo nào của \(M_p\).
  \item \textbf{Mốc bão hòa}: Nếu mốc \(t\) đã bị ép chọn bởi đúng \(K\) cuộc họp cố định, loại mốc \(t\) khỏi miền của tất cả cuộc họp còn lại.
\end{enumerate}

\begin{proposition}[Tính bảo toàn nghiệm của lọc miền B2B]
\label{prop:b2b-domain-preservation}
Thuật toán lọc miền dừng sau số bước hữu hạn, loại bỏ duy nhất các giá trị không thể thuộc bất kỳ lịch hợp lệ nào, bảo toàn trọn vẹn tập nghiệm khả thi, và chỉ thông báo vô nghiệm khi bài toán thực sự không có nghiệm.
\end{proposition}

\begin{proof}
Mỗi bước loại giá trị đều dựa trên các điều kiện cần về tính khả thi: vi phạm cận thứ tự tiên quyết, không có hỗ trợ trong ghép cặp đại biểu, hoặc vượt quá sức chứa địa điểm \(K\). Vì miền ban đầu là hữu hạn, quá trình xóa giá trị phải dừng sau tối đa \(\sum_{m}\card{D_m^0}\) bước.
\end{proof}

Phép biểu diễn \mbox{\textsc{Reduced}} áp dụng biến gán \(x_{m,t}\) chỉ trên miền mốc đã lọc \(A_m = D_m\), thay vì miền đầy đủ \(A_m = T\) của mô hình gốc \mbox{\textsc{ORG-MaxSAT}} \parencite{bofill2022b2b}.

\subsection{Biểu diễn thứ tự tiên quyết bằng hậu tố thưa chia sẻ (\textsc{SparseSuffix})}

Xét quan hệ tiên quyết \((i,j,d)\in\mathcal R_g\). Biểu diễn từng cặp (\mbox{\textnormal{\textsc{Pairwise}}}) thêm các mệnh đề nhị phân \(\neg x_{i,t}\lor\neg x_{j,u}\) với mọi \(t\in A_i, u\in A_j\) thỏa \(t+d>u\), có thể phát sinh \(O(\card{A_i}\card{A_j})\) mệnh đề cho một quan hệ.

Phép biểu diễn hậu tố thưa (\mbox{\textnormal{\textsc{SparseSuffix}}}) \parencite{nguyen2026compactsatmaxsatencodings} chỉ khởi tạo biến phụ tại các điểm cắt nội bộ thực sự xuất hiện trong miền tiền nhiệm. Với miền đã sắp thứ tự \(A_i = (a_{i,0},\ldots,a_{i,n_i-1})\), giả sử các điểm cắt xuất hiện là \(c_1 < c_2 < \cdots < c_s\). Đặt \(c_{s+1}=n_i\) và \(q_{i,n_i}=\bot\), biến đại diện hậu tố được định nghĩa quy nạp:
\begin{align}
  c(i,j,u) &= \min\set{r\in\set{0,\ldots,n_i-1} \suchthat a_{i,r} > u - d}, \notag\\
  q_{i,c_\ell} &\leftrightarrow \left(\bigvee_{r=c_\ell}^{c_{\ell+1}-1} x_{i,a_{i,r}}\right) \lor q_{i,c_{\ell+1}}, \label{eq:b2b-sparse-suffix}\\
  x_{j,u} &\rightarrow \neg q_{i,c(i,j,u)}. \label{eq:b2b-suffix-link}
\end{align}
Khi hai quan hệ tiên quyết cùng dẫn đến một điểm cắt \(c\) trên miền \(A_i\), biến hậu tố \(q_{i,c}\) được tái sử dụng hoàn toàn.

\begin{proposition}[Tương đương của SparseSuffix]
\label{prop:b2b-suffix-equivalence}
Phép biểu diễn \mbox{\textnormal{\textsc{SparseSuffix}}} và phép biểu diễn từng cặp \mbox{\textnormal{\textsc{Pairwise}}} mô tả chính xác cùng một tập phân gán cuộc họp--mốc thời gian.
\end{proposition}

\begin{proof}
Từ Công thức~\eqref{eq:b2b-sparse-suffix}, biến phụ \(q_{i,c}\) đúng khi và chỉ khi cuộc họp \(i\) được gán cho mốc có chỉ số ít nhất là \(c\). Mệnh đề kéo theo~\eqref{eq:b2b-suffix-link} do đó loại bỏ chính xác các cặp mốc \((t,u)\) vi phạm điều kiện \(t+d \le u\).
\end{proof}

\begin{figure}[tb]
\centering
\begin{tikzpicture}[satfig,x=8.5mm,y=7.5mm]
  %--- Diagram (a): Predecessor-domain cut ---
  \node[anchor=west,font=\rmfamily\bfseries\small] at (0,3.5) {(a) Điểm cắt miền tiền nhiệm};
  \node[anchor=east,font=\rmfamily\small] at (0.6,2.1) {$A_i$};
  
  \foreach \x/\v in {0.8/1, 1.8/3, 2.8/4, 3.8/7} {
    \draw[satbox] (\x,1.75) rectangle +(0.8,0.7);
    \node[font=\rmfamily\small] at (\x+0.4,2.1) {$\v$};
  }
  
  \draw[Gold,dashed,very thick] (2.7,1.5)--(2.7,2.7);
  \node[Gold,anchor=south,font=\rmfamily\bfseries\small] at (2.7,2.7) {$c=2$};
  
  \node[anchor=west,font=\rmfamily\small] at (4.8,2.1) {$x_{j,4}$};
  \node[anchor=west,font=\rmfamily\small] at (0.8,1.0) {$\neg x_{j,4} \lor \neg q_{i,2}$};
  \node[anchor=west,font=\rmfamily\small] at (0.8,0.3) {$q_{i,2} \leftrightarrow x_{i,4} \lor x_{i,7}$};

  %--- Separator ---
  \draw[Rule,thick] (6.5,0.1)--(6.5,3.6);

  %--- Diagram (b): Schedule span & idle-time threshold ---
  \node[anchor=west,font=\rmfamily\bfseries\small] at (7.2,3.5) {(b) Đánh giá ngưỡng chờ nội bộ};
  \node[anchor=east,font=\rmfamily\small] at (7.8,2.1) {mốc};
  
  \draw[satbox,fill=PaleBlue,draw=Blue,thick] (8.0,1.75) rectangle +(0.8,0.7);
  \node[font=\rmfamily\small] at (8.4,2.1) {$3$};
  
  \draw[satbox,fill=SoftGray] (9.0,1.75) rectangle +(0.8,0.7);
  \node[font=\rmfamily\small] at (9.4,2.1) {$4$};
  
  \draw[satbox,fill=PaleBlue,draw=Blue,thick] (10.0,1.75) rectangle +(0.8,0.7);
  \node[font=\rmfamily\small] at (10.4,2.1) {$5$};
  
  \draw[satbox,fill=SoftGray] (11.0,1.75) rectangle +(0.8,0.7);
  \node[font=\rmfamily\small] at (11.4,2.1) {$6$};
  
  \node[anchor=south,font=\rmfamily\small,text=Blue] at (8.4,2.55) {$f_{p,3}=1$};
  \node[anchor=south,font=\rmfamily\small,text=Blue] at (10.4,2.55) {$R_{p,5}=1$};
  
  \draw[satdim] (8.4,1.4)--(10.4,1.4);
  \node[anchor=north,font=\rmfamily\small] at (9.4,1.35) {$\card{M_p}=2 \implies I_p \ge 1$};
\end{tikzpicture}
\caption{Kỹ thuật biểu diễn trong B2B: (a) Biến hậu tố thưa \mbox{\textnormal{\textsc{SparseSuffix}}} tại điểm cắt $c=2$; (b) Xác định ngưỡng thời gian chờ nội bộ từ mốc bắt đầu và mốc kết thúc.}
\label{fig:b2b-encoding-ideas}
\end{figure}

\subsection{Mã hóa ngưỡng thời gian chờ từ nhịp hoạt động}

Để xác định \(I_p(S)\), gọi \(y_{p,t}\leftrightarrow\bigvee_{m\in M_p: t\in A_m}x_{m,t}\) là biến chỉ báo đại biểu \(p\) bận tại mốc \(t\). Định nghĩa biến tích lũy tiền tố \(P_{p,t}\leftrightarrow P_{p,t-1}\lor y_{p,t}\) và hậu tố \(R_{p,t}\leftrightarrow y_{p,t}\lor R_{p,t+1}\). Mốc bắt đầu đầu tiên của đại biểu \(p\) được nhận diện bằng:
\[
  f_{p,t} \leftrightarrow y_{p,t} \land \neg P_{p,t-1}.
\]
Gọi \(\theta_{p,k}\) là biến mệnh đề đại diện cho điều kiện \(I_p(S)\ge k\). Mối liên hệ giữa mốc đầu và mốc cuối được biểu diễn chính xác bằng:
\begin{equation}
  f_{p,t} \rightarrow \left(\theta_{p,k} \leftrightarrow R_{p, t + \card{M_p} + k - 1}\right)
  \qquad \forall t\in T, 1\le k\le U_p,
  \label{eq:b2b-span-threshold}
\end{equation}
trong đó \(U_p\) là cận trên hợp lệ của thời gian chờ đối với đại biểu \(p\).

Cuối cùng, với \(U=\max_{p\in P^\star} U_p\) và \(1\le k\le U\), biểu diễn định nghĩa biến ngưỡng cực đại \(\Theta_k^{\max}\leftrightarrow\bigvee_{p\in P^\star}\theta_{p,k}\), cực tiểu \(\Theta_k^{\min}\leftrightarrow\bigwedge_{p\in P^\star}\theta_{p,k}\), và biến mức chênh lệch:
\begin{equation}
  g_k \leftrightarrow \Theta_k^{\max} \land \neg \Theta_k^{\min}, \qquad \text{với } \sum_{k=1}^U g_k = \Delta_I(S).
  \label{eq:b2b-range-sum}
\end{equation}

\subsection{Phương pháp giải và kết quả thực nghiệm}

Phép biểu diễn B2B thu gọn được đánh giá thực nghiệm rộng rãi trên 126 bộ dữ liệu chuẩn chính thức (\emph{All-126}) và 100 bộ dữ liệu biến đổi mật độ thứ tự (\emph{Stress-100}) \parencite{nguyen2026compactsatmaxsatencodings}. Thử nghiệm so sánh ba phương pháp tối ưu dựa trên công thức SAT/MaxSAT (UWrMaxSAT, \textsc{Non-Incremental SAT}, và \textsc{Incremental SAT} sử dụng Totalizer kết hợp giả định assumption) với các công cụ thương mại hàng đầu đại diện cho MIP (Gurobi, CPLEX) và CP (CP Optimizer).

\begin{resultbox}{Kết quả tiêu biểu cho B2B Meeting Scheduling}
Trên tập dữ liệu chuẩn \emph{All-126}:
\begin{itemize}
  \item Phép biểu diễn thu gọn giảm trung vị \textbf{40.3\%} số mệnh đề, \textbf{55.9\%} dung lượng bộ nhớ đỉnh và \textbf{14.0\%} tổng thời gian chạy so với mô hình \mbox{\textsc{ORG-MaxSAT}} công bố \parencite{bofill2022b2b}.
  \item Riêng bước lọc miền (\mbox{\textsc{Reduced}}) đóng góp giảm \textbf{24.1\%} số biến gán và \textbf{16.2\%} số mệnh đề.
  \item Biểu diễn hậu tố \mbox{\textnormal{\textsc{SparseSuffix}}} kết hợp với bao đóng bắc cầu \mbox{\textnormal{\textsc{Closure-}}\ensuremath{E^\star}} giúp giảm số mệnh đề đến \textbf{5.5\%} ở các mức mật độ thứ tự cao.
  \item Cả ba phương pháp SAT và MaxSAT đều giải thành công 100\% các bài toán chính thức với thời gian trung vị vượt trội Gurobi MIP (\(0.521\) giây); trong đó \mbox{\textnormal{\textsc{Incremental SAT}}} đạt thời gian trung vị nhanh nhất (\textbf{0.227} giây).
\end{itemize}
\end{resultbox}

\section{Peak power: công suất đỉnh trong cân bằng dây chuyền}

\subsection{Mô hình đỉnh công suất}

Trong SALBP-3PM, tác vụ \(i\) có thời lượng \(t_i\), công suất \(w_i\), được gán
vào một trạm và bắt đầu trong chu kỳ dài \(c\). Nếu \(A_{i\tau}\) biểu diễn
``tác vụ \(i\) đang hoạt động tại \(\tau\)'', tổng công suất là
\[
  W(\tau)=\sum_i w_iA_{i\tau},
  \qquad
  P=\max_{\tau=0}^{c-1}W(\tau).
\]
Mục tiêu là tối thiểu hóa \(P\).

Các biến trực tiếp cho quyết định gán và bắt đầu thuận tiện cho phép giải mã:
\[
  X_{ik}=1 \iff i\text{ ở trạm }k,
  \qquad
  S_{it}=1 \iff i\text{ bắt đầu tại }t.
\]
Các biến thứ tự cung cấp lớp lan truyền:
\[
  R_{ik}=1 \iff \text{trạm}(i)\le k,
  \qquad
  T_{it}=1 \iff \text{bắt đầu}(i)\le t.
\]
Giá trị chính xác xuất hiện tại điểm chuyển của chuỗi thứ tự. Nhờ vậy, ràng
buộc AMO về phép gán giảm từ dạng từng cặp bậc hai xuống một chuỗi tuyến tính;
quan hệ trước--sau giữa trạm hoặc giữa mốc thời gian cũng được viết bằng quan
hệ kéo theo giữa các ngưỡng.

\subsection{Bốn giao diện tối ưu}

Cùng một phép biểu diễn tính khả thi có thể kết nối với hàm mục tiêu theo bốn cách:
\begin{center}
\captionof{table}{Bốn giao diện tối ưu cho công suất đỉnh.}
\label{tab:peak-interfaces}
\begin{tabularx}{\textwidth}{@{}>{\bfseries\sffamily}p{.23\textwidth}YY@{}}
\toprule
Cơ chế & Cách thay cận & Thành phần cần giữ cố định\\
\midrule
Chặn bằng mệnh đề & Cấm mọi đỉnh không tốt hơn nghiệm đương nhiệm &
Ngữ nghĩa của biến đỉnh\\
Tinh chỉnh PB & Siết \(\sum_iw_iA_{i\tau}\le B\) tại mọi \(\tau\) &
Phép biểu diễn giả Boolean\\
\MaxSAT có trọng số & Đưa mức đỉnh vào mệnh đề mềm &
Ánh xạ chi phí với giá trị \(P\)\\
Biến chỉ báo thứ tự & Kích hoạt cận bằng các giả định &
Mệnh đề cơ sở và chuỗi thứ tự\\
\bottomrule
\end{tabularx}
\end{center}

Biến chỉ báo thứ tự phù hợp với giải tăng dần: bộ giải được dựng một lần ở cận
trên, sau mỗi nghiệm chỉ thay các giả định hoặc thêm mệnh đề chặn. Trong nghiên
cứu của \textcite{kieu2025power}, phép biểu diễn tuần tự gọn kết hợp giải tăng dần
chứng nhận được nhiều bài toán hơn cấu hình đối chiếu trên tập 89 bài kiểm
chuẩn; lợi ích rõ nhất ở nhóm khó.

\begin{keyidea}{Tách tính khả thi khỏi tối ưu}
Mô-đun khả thi mô tả lịch hợp lệ qua một giao diện độc lập với tìm kiếm tuyến
tính, tìm kiếm nhị phân hay \MaxSAT. Lớp hàm mục tiêu đọc các biến đã có và
cung cấp giao diện cận \(B\). Cách tách này làm cho chứng minh tính đúng ngắn
hơn và thí nghiệm phân tích thành phần (ablation study) rõ hơn.
\end{keyidea}

\section{Điều chỉnh lịch tàu bằng MaxSAT--DDD}

\subsection{Tại sao cần rời rạc hóa động}

Với tàu \(i\) và tài nguyên \(r\) trên tuyến cố định, đặt \(t^{i,r}\) là thời
điểm tàu đi vào tài nguyên. Quan hệ trước--sau trên đường đi có dạng
\[
  t^{i,q}\ge t^{i,r}+\ell_i^r,
\]
trong khi hai tàu chia sẻ tài nguyên phải chọn một trong hai thứ tự an toàn.
Nếu dựng toàn bộ miền thời gian đến khung thời gian \(H\), số biến theo mốc thời gian có thể
rất lớn.

\emph{Dynamic Discretization Discovery} (khám phá rời rạc hóa động, DDD) khởi
đầu bằng tập mốc thưa \(\Lambda^{i,r}\). Mô hình hạn chế cho một cận dưới.
Nghiệm được giải mã trên trục thời gian đầy đủ; nếu nghiệm vi phạm thời gian di
chuyển hoặc khoảng cách an toàn, chỉ các mốc liên quan mới được thêm.

DDD không bắt nguồn từ SAT. Thuật toán được phát triển cho thiết kế mạng dịch
vụ thời gian liên tục bằng cách giải lặp các mạng mở rộng thời gian một phần
và tinh chỉnh những mốc cần thiết \parencite{boland2017ddd}. Khi chuyển ý tưởng
sang MaxSAT, điều phải bảo toàn là \emph{tính hợp lệ của cận dưới} trong mô
hình hạn chế và quy tắc tinh chỉnh hữu hạn; bộ giải bên trong có thể thay đổi.
Đây là một ví dụ tiêu biểu cho quá trình tích hợp kiến trúc từ Nghiên cứu vận trù
vào SAT, thay vì chỉ thay phép biểu diễn của một ràng buộc.

\subsection{Vòng lặp}

\begin{enumerate}
  \item Giải mô hình \MaxSAT hạn chế trên phép rời rạc hóa hiện tại.
  \item Giải mã lịch và kiểm tra mọi quan hệ trước--sau trên đường đi, xung
        đột tài nguyên và giá trị độ trễ.
  \item Nếu có vi phạm, thêm các mốc cần thiết và lặp lại.
  \item Nếu lịch khả thi và chi phí giải mã bằng giá trị tối ưu của mô hình
        hạn chế, cận dưới gặp cận trên; nghiệm là tối ưu.
\end{enumerate}

\begin{algorithm}[H]
\caption{Vòng lặp MaxSAT--DDD có kiểm chứng}
\label{alg:maxsat-ddd}
\begin{minipage}{.94\textwidth}\small
\textbf{Đầu vào:} tập mốc hữu hạn đầy đủ \(\Lambda^\star\), phép rời rạc hóa ban đầu
\(\Lambda^0\subseteq\Lambda^\star\), bộ giải mã–kiểm tra và hàm chi phí.\\
\textbf{Đầu ra:} lịch tối ưu hoặc chứng nhận vô nghiệm trong miền đầy đủ.
\begin{enumerate}
  \item Với vòng \(r\), sinh hoặc cập nhật mô hình hạn chế trên
  \(\Lambda^r\), giải tối ưu và ghi cận dưới \(L_r\).
  \item Nếu mô hình hạn chế vô nghiệm và là một mô hình nới lỏng hợp lệ của mô hình
  đầy đủ, kết luận mô hình đầy đủ vô nghiệm.
  \item Giải mã nghiệm. Bộ kiểm tra xác nhận một lịch đầy đủ có chi phí \(U_r\) nếu
  khả thi; đồng thời, khi lịch chưa khả thi hoặc \(L_r<U_r\), thuật toán phải trả về một
  bằng chứng \(C_r\) chỉ ra phần rời rạc hóa đang thiếu cho feasibility hoặc
  hàm mục tiêu.
  \item Nếu có lịch khả thi và \(U_r=L_r\), trả lịch cùng hai cận bằng nhau.
  \item Nếu chưa hội tụ, bước tinh chỉnh phải chọn ít nhất một mốc trong
  \(\Lambda^\star\setminus\Lambda^r\) được bằng chứng \(C_r\) yêu cầu; đặt
  \(\Lambda^{r+1}\supsetneq\Lambda^r\) và lặp.
\end{enumerate}
\end{minipage}
\end{algorithm}

\begin{proposition}[Hội tụ hữu hạn của DDD]
\label{prop:ddd-finite}
Giả sử (i) tập mốc đầy đủ \(\Lambda^\star\) hữu hạn; (ii) mỗi mô hình hạn chế
cho một cận dưới hợp lệ; (iii) bộ kiểm tra đúng; và (iv) mỗi vòng chưa kết thúc
thêm ít nhất một mốc mới từ \(\Lambda^\star\). Khi đó
Thuật toán~\ref{alg:maxsat-ddd} kết thúc sau không quá
\(\card{\Lambda^\star\setminus\Lambda^0}+1\) lần giải mô hình hạn chế và chỉ
trả nghiệm tối ưu hoặc kết luận vô nghiệm đúng.
\end{proposition}

\begin{proof}
Do chuỗi
\(\Lambda^0\subsetneq\Lambda^1\subsetneq\cdots\subseteq\Lambda^\star\)
tăng nghiêm ngặt trong một tập hữu hạn, số lần tinh chỉnh bị chặn bởi
\(\card{\Lambda^\star\setminus\Lambda^0}\). Nếu thuật toán dừng với
\(U_r=L_r\), \(L_r\) là cận dưới cho giá trị tối ưu đầy đủ còn \(U_r\) là chi
phí của một nghiệm đầy đủ được bộ kiểm tra xác nhận, nên cả hai bằng giá trị
tối ưu. Nếu đến \(\Lambda^\star\), mô hình hạn chế chính là mô hình đầy đủ; bộ giải quyết định
tối ưu hoặc vô nghiệm. Kết luận vô nghiệm sớm chỉ hợp lệ dưới giả thiết
nới lỏng ở (ii).
\end{proof}

\paragraph{Biên hội tụ và chiến lược tinh chỉnh.}
Ở biên xấu nhất, DDD khôi phục toàn bộ lưới \(\Lambda^\star\). Mỗi bước tinh
chỉnh vì thế phải thêm một mốc mới được bộ kiểm tra suy ra từ vi phạm cụ thể.
Với miền thời gian liên tục, một đối chứng thời gian liên tục hoặc một định lý
về tập ứng viên hữu hạn cung cấp điều kiện dừng. Ba thành phần này---mốc mới, bộ kiểm tra và
tập mốc đầy đủ---xác định đồng thời tính đúng và tốc độ hội tụ.

\begin{figure}[tb]
\centering
\begin{tikzpicture}[satfig,x=10mm,y=8mm]
  \draw[->,Rule] (0,0)--(10,0) node[right,satlabel]{thời gian};
  \foreach \t in {0,2,5,8,10}{
    \draw[Rule] (\t,.12)--(\t,-.12);
    \node[below,satlabel] at (\t,-.15) {\(\t\)};
  }
  \node[left,satlabel] at (0,1.1) {vòng \(r\)};
  \foreach \t in {0,5,10}
    \node[satstate,minimum size=6mm] at (\t,1.1) {};
  \draw[Blue,thick] (0,1.1)--(5,1.1)--(10,1.1);
  \draw[Gold,very thick] (5,1.1)--(8,1.1)
    node[midway,above,satlabel,text=Gold]{vi phạm được giải mã};
  \node[left,satlabel] at (0,2.5) {vòng \(r+1\)};
  \foreach \t in {0,5,10}
    \node[satstate,minimum size=6mm] at (\t,2.5) {};
  \node[satstate,satwarn,minimum size=6mm] at (8,2.5) {};
  \draw[Teal,thick] (0,2.5)--(5,2.5)--(8,2.5)--(10,2.5);
  \draw[satedge] (8,1.55)--(8,2.1)
    node[midway,right,satlabel]{thêm mốc \(8\)};
\end{tikzpicture}
\caption{DDD chỉ thêm mốc tại nơi nghiệm của mô hình hạn chế thiếu thông tin.
Lưới không được làm dày đồng loạt trên toàn bộ khung thời gian.}
\label{fig:ddd-refinement}
\end{figure}

Lan truyền quan hệ trước--sau có vai trò thu hẹp miền trước khi sinh \CNF. Nếu một
lựa chọn thứ tự \(i\prec j\) kéo theo một chuỗi thời điểm tối thiểu, các mốc
không đạt được bị loại ngay. AMO kết hợp tiếp tục chọn phép biểu diễn khác nhau
theo kích thước miền: từng cặp cho tập rất nhỏ, tuần tự hoặc chỉ huy cho tập lớn hơn.

\begin{resultbox}{Điểm rút ra từ Train MaxSAT--DDD}
Trong nghiên cứu của \textcite{kieu2026train}, hiệu quả thay đổi theo
dạng hàm độ trễ: \MaxSAT nổi bật với chi phí bậc thang và làm tròn, còn MIP cạnh
tranh hơn với chi phí liên tục. Giá trị thiết kế của DDD nằm ở khả năng chỉ mở
rộng miền tại nơi nghiệm hạn chế thực sự thiếu thông tin.
\end{resultbox}

\section{Chọn SAT, SMT, CP hay MIP}

Không có bộ giải thống trị mọi cấu trúc. Một quy tắc thực dụng là nhìn vào phần
khó nhất của mô hình:
\begin{itemize}
  \item Chọn \SAT/\MaxSAT khi lôgic rời rạc, cửa sổ và xung đột lặp lại; có lợi
        từ học mệnh đề và giải tăng dần.
  \item Chọn SMT khi bất đẳng thức thời gian nguyên là lõi, còn phần Boolean
        chỉ quyết định thứ tự hoặc chế độ \parencite{bofill2020rcpsp}.
  \item Chọn CP khi tích lũy/không chồng lấn là ràng buộc toàn cục tự nhiên và
        lan truyền miền có thể khai thác trực tiếp.
  \item Chọn MIP khi hàm mục tiêu tuyến tính mạnh và mô hình nới lỏng cung cấp cận tốt.
\end{itemize}

Phương pháp kết hợp nên được thiết kế như một phân chia nhiệm vụ suy diễn: mô-đun
nào tạo cận, mô-đun nào xử lý phép tuyển và bằng cách nào hai phía trao đổi
thông tin. Việc chỉ gọi nhiều bộ giải trên cùng một mô hình không tự tạo ra
một kiến trúc kết hợp.

\subsection{Một cây quyết định theo cấu trúc}

\begin{enumerate}
  \item Nếu khung thời gian nhỏ và mỗi tác vụ có ít giá trị bắt đầu, SAT/MaxSAT theo
        mốc thời gian là đối chứng minh bạch.
  \item Nếu khung thời gian lớn nhưng các ràng buộc trước--sau/sai phân dày, thử mã
        hóa thứ tự, SMT hoặc LCG trước khi mở toàn bộ trục thời gian.
  \item Nếu tài nguyên tích lũy là nút thắt và có suy luận năng lượng mạnh,
        CP/LCG có ưu thế cấu trúc mà xung đột từng cặp khó thay thế.
  \item Nếu hàm mục tiêu/mô hình nới lỏng tuyến tính tạo cận mạnh, MIP là đối
        chứng bắt buộc; SAT chỉ có lợi khi phép tuyển Boolean hoặc đối xứng chi phối.
  \item Nếu chỉ một phần nhỏ mốc thời gian quan trọng, xét DDD/tinh chỉnh và
        chứng minh rõ tính hợp lệ của cận dưới.
\end{enumerate}

\section{Bài học thiết kế}

\begin{summarybox}
\begin{enumerate}
  \item Phân biệt biến thời điểm, biến hoạt động và biến thứ tự; viết phép
        liên kết trước khi tối ưu mệnh đề.
  \item Tính không ngắt là một ràng buộc khối, không phải chỉ là ràng buộc đếm.
  \item Với các cửa sổ chồng lấn, hiệu quả đến từ kỹ thuật chia sẻ ngưỡng hoặc
        bộ đếm giữa các cửa sổ.
  \item Hàm mục tiêu nên giao tiếp với mô-đun khả thi qua một giao diện cận.
  \item Rời rạc hóa động phù hợp khi miền lớn nhưng nghiệm chỉ cần một số
        ít mốc quan trọng.
\end{enumerate}
\end{summarybox}

\subsection*{Câu hỏi nghiên cứu}
\begin{enumerate}
  \item Với cùng ngữ nghĩa khối, hãy so sánh biểu diễn bằng biến bắt đầu và biểu diễn
        bằng điểm chuyển về số biến, số mệnh đề và khả năng lan truyền đơn vị.
  \item Thiết kế một Shared Counter cho hai ALSC có cửa sổ chồng lấn nhưng cận
        \(k\) khác nhau.
  \item Viết điều kiện dừng chính xác cho một vòng DDD khi chi phí hạn chế là
        cận dưới và lịch giải mã cho cận trên.
\end{enumerate}
